\chapter{Transformador}
\label{cap:Transformers}
La aparición de los transformadores (\textit{transformer}) ha marcado un antes y un después en el campo del aprendizaje profundo (\textit{deep learning}). Introducida a finales de 2017 por un equipo de investigadores de Google  en el célebre artículo ``\textit{Attention is all you need}'' \cite{vaswani2017attention_arxiv}, esta arquitectura desafió el paradigma dominante al prescindir por completo de las clásicas redes neuronales recurrentes (RNN) para el procesamiento de secuencias.

Su gran innovación radica en el mecanismo de autoatención (\textit{self-attention}), el cual permite determinar qué partes de la información de entrada son más relevantes en cada momento. El nombre de transformador deriva precisamente de su función: transformar un conjunto de vectores iniciales en un nuevo conjunto de vectores de la misma dimensión. Sin embargo, esta nueva representación de la entrada es mucho más rica en contexto, significado e interpretación de los datos iniciales, lo que resulta fundamental para resolver con éxito tareas posteriores.

Las implicaciones que ha tenido esta nueva arquitectura han sido sorprendentes, pues aunque nacieron con el objetivo de revolucionar el Procesamiento de Lenguaje Natural (NLP), su impacto ha trascendido esta disciplina y  forman la base de modelos de visión artificial y sistemas multimodales. Tal ha sido la magnitud de este hito tecnológico que los  modelos extensos de lenguaje basan su diseño en esta arquitectura: la familia de modelos GPT (\textit{Generative Pre-Trained Transformer}) de OpenAI, el modelo BERT (\textit{Bidirectional Encoder Representations from Transformers}) y modelos como Gemini de Google entre otros.

Este capítulo pretende asentar y desgranar en detalle la arquitectura interna del transformador, analizando el funcionamiento de todos los ingredientes que lo componen y, sobre todo, el mecanismo de autoatención que hace posible el funcionamiento de este.

\section{Procesamiento}
\label{procesamiento}
La entrada de un transformador será un conjunto de vectores $\{\x_n\}$ de dimensión $D_i$ donde $n = 1, \dots,N$. Dichos vectores recibirán el nombre de \textit{tokens}, donde un \textit{token} podrá representar una palabra dentro de una frase, partes de una palabra o incluso de una imagen. Además, dado un \textit{token} $\x_n^T = ( \x_{n1}, \x_{n2}, \dots,\x_{nD_i})$ llamaremos a sus componentes $\x_{nm}$ rasgos.

Los vectores o \textit{tokens} $\{\x_n\}$ los agruparemos en una matriz $\mathbf{X}$ de dimensión $N \times D_i$ donde la fila $n$-ésima representará el  \textit{token} ${\x_n ^{T}}$ (ver \ref{fig:matriz_X}).

\begin{figure}[t]
    \vspace{0.8cm}
    \centering
    \MatrizX
    \caption{Esquema de la matriz de datos.}
    \label{fig:matriz_X}
    \vspace{0.8cm}
\end{figure}
El transformador será una función no lineal que mapea la matriz de 
entrada $\mathbf{X}$ hacia una nueva representación $\mathbf{Y}$ de la
misma dimensión $N \times D_i$:
\begin{equation}
    \mathbf{Y} = \text{Transformador}[\mathbf{X}]
\end{equation}
La salida $\mathbf{Y}$ será  capaz de capturar la relación que tienen los 
vectores de entrada gracias a un mecanismo de autoatención. 
Esto permitirá que cada vector resultante $\y_i^T$ de $\mathbf{Y}$ integre 
el contexto global de la secuencia, capturando la información relevante de 
todos los demás elementos.

Supongamos que nuestro conjunto inicial de \textit{tokens} es $\x_1,\x_2, \dots,\x_N$ y queremos convertirlo en  $\y_1,\y_2, \dots,\y_N$, de tal manera que cada uno de los $\y_n$ no solo dependa del vector de entrada $\x_n$ sino también de todos los demás vectores $\x_m$.
La manera más elegante y sencilla de expresar esta dependencia es mediante una combinación lineal:
\begin{equation}
    \y_n = \sum_{m=1}^{N} a_{nm}\x_m \qquad n = 1,2,\dots, N
    \label{combinacion_lineal}
\end{equation}
donde los coeficientes $a_{nm}$ recibirán el nombre de \textbf{coeficientes de atención}.
Estos coeficientes $a_{nm}$ tienen como objetivo representar la influencia de los vectores $\x_m$ sobre el vector de salida $\y_n$. Valores muy cercanos al cero indicarán poca influencia y valores más grandes indicarán una mayor influencia sobre aquel. Sin embargo, también queremos asegurar el hecho de que si una entrada particular ``pone más atención''  sobre el vector salida, entonces el resto de entradas deberán ``poner menos atención''  sobre esta. Matemáticamente:

\begin{equation}
    a_{nm} \geq 0 \qquad n,m = 1,2,\dots, N
    \label{atencion_positiva}
\end{equation}

\begin{equation}
    \sum_{m=1}^{N} a_{nm} = 1 \qquad n=1,2,\dots, N
    \label{atencion_suma_1}
\end{equation}

\section{Coeficientes de atención}
El objetivo es ahora determinar cómo calcular dichos 
\textbf{coeficientes de atención}. Iremos de menos a más, comenzando con 
un modelo sencillo al que, progresivamente, iremos añadiendo complejidad y 
detalles.

Para ponernos en situación, supongamos un grupo de filósofos sentados 
alrededor de una mesa redonda debatiendo y colaborando en la creación de una 
nueva teoría filosófica. Cada filósofo posee una etiqueta que identifica su 
corriente de pensamiento o especialidad, la cual actúa como \textbf{clave}, 
y un conocimiento profundo que puede aportar al grupo, que actúa como 
\textbf{valor}. Cuando un filósofo necesita inspiración o tiene dudas sobre 
algo formula una pregunta al resto de la mesa, la cual actúa como 
\textbf{consulta}. De todos los integrantes que hay en la mesa, el filósofo 
que lanza la pregunta necesita saber a quién debe ``prestar más atención'', 
es decir, qué compañero tiene la especialidad que mejor encaje con su 
pregunta. En otras palabras, quiere encontrar a aquel filósofo cuya clave 
coincida mejor con su consulta. Sin embargo, como buen filósofo, no solo 
quiere tener en cuenta la respuesta de uno sino la de todos los filósofos 
en la mesa. Para ello, el filósofo comparará su consulta con todas las 
claves de sus compañeros e incluso con su propia clave (de donde deriva el 
término de autoatención) y asignará una distribución de pesos que medirán 
cuánta atención prestar a cada filósofo. Finalmente, cada filósofo le dará 
su respuesta, es decir, su valor, y el filósofo que lanza la pregunta 
considerará y ponderará estas respuestas según los pesos de atención. De 
esta manera, el filósofo tendrá en cuenta las respuestas de todos los 
participantes pero a algunos les hará más caso que a otros. Cabe destacar 
que este proceso es simultáneo, pues cada uno de los integrantes lanzará su 
propia pregunta al resto de filósofos, actualizando y enriqueciendo su 
conocimiento al mismo tiempo.

En una primera aproximación simplificada, supondremos $N$ filósofos en la mesa.
Cada uno de ellos será representado mediante un vector $\x_m$.
Las consultas y las claves serán representados 
mediante los vectores $\mathbf{q}_m$ y $\mathbf{k}_m$  respectivamante, 
mientras que el valor
será el propio filósofo, es decir, el vector $\x_m $
Dado un filósofo $\x_n$, será necesario comparar su correspondiente consulta
$\mathbf{q}_n$ con las claves $\mathbf{k}_m$ de
los filósofos presentes en la mesa. Para realizar la comparación y medir el 
grado de similitud emplearemos el producto escalar, por lo que podríamos 
definir, inicialmente, los coeficientes de atención como:
\begin{equation}
    a_{nm} = \mathbf{q}_n^{{T}}  \cdot \mathbf{k}_m \qquad \text{(consulta} \cdot \text{clave)} 
\end{equation}
donde  $\cdot$ denota el producto escalar. 
Sin embargo, al no cumplir \eqref{atencion_positiva} ni 
\eqref{atencion_suma_1}, es necesario aplicar la función Softmax 
sobre el conjunto de productos escalares obtenidos. Se definen entonces los coeficientes de atención 
como:
\begin{equation}
    a_{nm} := \text{Softmax}_m(\mathbf{q}_n^{T} \mathbf{k}_m) = 
    \frac{\exp(\mathbf{q}_n^{T} \mathbf{k}_m)}{\sum_{m'=1}^{N} \exp(\mathbf{q}_n^{T} \mathbf{k}_{m'})} 
    \qquad  m = 1, \dots, N
    \label{softmax_anm_componente}
\end{equation}
donde cada uno de estos coeficientes representaría la atención que pone 
el filósofo $\x_n$, el  que realiza la consulta, sobre el filósofo $\x_m$, 
el que 
puede responder a su consulta.

Según comentábamos, el proceso es simultáneo y cada vector 
$\x_n$ de $\mathbf{X}$, mediante \eqref{combinacion_lineal}, genera su 
propio vector $\y_n$, que, conceptualmente y aplicado a la analogía,
representa la conclusión a la que llega el filósofo $\x_n$.

Todas estas ecuaciones, \eqref{combinacion_lineal} y 
\eqref{softmax_anm_componente}, se pueden expresar matricialmente como:
\begin{equation}
    \mathbf{Y} = \text{Softmax}[\mathbf{Q}  \mathbf{K}^T]\mathbf{X}
    \label{eq:softmax_anm_matriz}
\end{equation}
donde $\mathbf{Y}$ tiene la  misma dimensión $N\times D_i$ que $\mathbf{X}$ y 
cuyas filas son los vectores de salida $\y_1^T,\y_2^T,\dots,\y_m^T$.

Si observamos la ecuación \eqref{eq:softmax_anm_matriz}, se 
aprecia que la transformación de $\{\mathbf{x}_m\}$ a $\{\mathbf{y}_m\}$ es 
rígida, ya que hemos asumido que las consultas y las claves son valores fijos 
que no dependen de la entrada $\mathbf{X}$. Bajo esta premisa, cada filósofo 
aporta siempre el mismo valor independientemente del filósofo, lo que 
otorga al modelo muy poca flexibilidad para decidir, en función de la 
situación, a quién debe prestar más atención.  
Es por ello que nos gustaría que el modelo tuviese capacidad de aprendizaje, 
es decir, que no se limite a una transformación fija, sino que aprenda a 
ajustar las representaciones de las consultas, las claves y los valores 
durante su entrenamiento. Para ello introduciremos hiperparámetros en el cálculo
de las consultas, las claves y los valores.

\section{Consulta, Clave y Valor}

Como comentábamos al final de la sección anterior, la ecuación \eqref{eq:softmax_anm_matriz} carece de parámetros y no hay capacidad de aprendizaje. Nos interesará entonces transformar la entrada $\mathbf{X}$ mediante alguna transformación lineal:
\begin{equation}
    \mathbf{\Tilde{X}} = \mathbf{X}\mathbf{\Omega} + \boldsymbol{\beta}
    \label{transformacion_X_matriz_pesos}
\end{equation}
donde $\mathbf{\Omega}$ es una matriz de pesos de dimensión $D_i \times D_i$ 
y $\boldsymbol{\beta}$ son los sesgos, como si de una red neuronal de una 
sola capa se tratase\footnote{
La ecuación \eqref{transformacion_X_matriz_pesos} se puede expresar 
matricialmente como 
$\mathbf{\Tilde{X}}= \begin{bmatrix} \mathbf{X} & 1 \end{bmatrix} \begin{bmatrix}
    \mathbf{\Omega} \\ \boldsymbol{\beta}
\end{bmatrix}$.
De aquí en adelante y con el fin de aligerar la notación  asumiremos que 
los sesgos están implícitos en la propia matriz de  pesos 
$\mathbf{\Omega}$ y \eqref{transformacion_X_matriz_pesos} equilvadrá a 
$ \mathbf{\Tilde{X}} = \mathbf{X}\mathbf{\Omega}$.}.
De esta manera, nos interesaría que el modelo fuese capaz de aprender características independientes para la consulta; distintas características para las claves y, también, otras características para los valores. En consecuencia, definiremos 3 matrices de pesos independientes con su correspondiente transformación lineal\footnote{
Se emplean las letras \textbf{Q}, \textbf{K} y \textbf{V} por su terminología en inglés: \textit{Query} (consulta), \textit{Key} (clave) y \textit{Value} (valor).
}:
\begin{equation}
    \mathbf{Q} = \mathbf{X}\mathbf{\Omega}_q 
    \label{omega_q}
\end{equation}
\begin{equation}
    \mathbf{K} = \mathbf{X}\mathbf{\Omega}_k 
    \label{omega_k}
\end{equation}
\begin{equation}
    \mathbf{V} = \mathbf{X}\mathbf{\Omega}_v
    \label{omega_v}
\end{equation}

Los parámetros de las matrices de pesos $\mathbf{\Omega}_q$, $\mathbf{\Omega}_k$ y $\mathbf{\Omega}_v$ irán ajustándose durante el entrenamiento de manera independiente, consiguiendo así esa flexibilidad que buscábamos. De esta manera, juntando \eqref{omega_q}, \eqref{omega_k} y \eqref{omega_v} con \eqref{eq:softmax_anm_matriz} obtenemos :
\begin{equation}
    \mathbf{Y} = \text{Softmax}[\mathbf{Q}  \mathbf{K}^T]\mathbf{V}
    \label{eq:eq_transformer}
\end{equation}

Las matrices $\mathbf{\Omega}_q $, $\mathbf{\Omega}_k $ y $\mathbf{\Omega}_v$ tendrán dimension $D_i \times D_i$ con el fin de garantizar la correcta multiplicación de las matrices \textbf{Q} y \textbf{K}, y asegurar que la salida \textbf{Y} tenga la misma dimension $N \times D_i$ que \textbf{X}.

Este mecanismo  recibe el nombre, en inglés, de \textit{dot-product self-attention} o de \textit{attention-head}. Nosotros lo llamaremos \textit{capa de autoatención} o \textit{mecanismo de autoatención} y lo denotaremos por:
\begin{equation}
    \mathbf{Y}[\mathbf{X}] = \textbf{Atención}[\mathbf{Q}, \mathbf{K}, \mathbf{V}]
    \label{atencion}
\end{equation}

En una red neuronal convencional cada entrada  se multiplica por un peso fijo, de modo que si la red decide ignorar una entrada, lo hace para todos los datos. Por el contrario, con este nuevo mecanismo, cada activación se multiplica por el coeficiente de atención \textbf{dependiente} de los datos de entrada. Esta distinción permite al modelo prestar más atención a unas entradas que a otras según el contexto (entradas).



\section{Atención Múltiple}

\begin{figure}[t]
    \vspace{0.8cm}
    \centering
    \ConcatenarMatrices
    \caption{Concatenacion de capas de atención}
    \label{fig:concatenar_matrices}
    \vspace{0.8cm}
\end{figure}

En ocasiones, suele haber múltiples mecanismos o patrones distintos a los que hay que prestar atención en un mismo problema. Por ejemplo, en el lenguaje natural algunos patrones suelen estar asociados al tiempo verbal, como el sufijo ``-ed'' del inglés que hace referencia al pasado; mientras que otros patrones tienen relación con la semántica o el significado de las palabras. Por ello, resulta interesante tener varias capas de autoatención trabajando simultáneamente y cada una de ellas  con sus propias matrices de consulta, clave y valor.

Supongamos que tenemos $H$ capas de atención independientes de dimensión $N \times D_i$:
\begin{equation}
    \mathbf{A}_h[\mathbf{X}]= \textbf{Atención}[\mathbf{Q}_h,\mathbf{K}_h, \mathbf{V}_h]
\end{equation}
donde las matrices de consulta, clave y valor son independientes, es decir:
\begin{equation}
    \mathbf{Q}_h = \mathbf{X}\mathbf{\Omega}_{q}^h 
    \label{omega_q_h}
\end{equation}
\begin{equation}
    \mathbf{K}_h = \mathbf{X}\mathbf{\Omega}_{k}^h
    \label{omega_k_h}
\end{equation}
\begin{equation}
    \mathbf{V}_h = \mathbf{X}\mathbf{\Omega}_{v}^h
    \label{omega_v_h}
\end{equation}
para cada $h = 1,2,\dots,H$. Con todo ello, concatenaremos las capas 
$\mathbf{A}_h$ y aplicaremos una transformación lineal mediante una matriz $\mathbf{\Omega_o}$ cuyos elementos se irán ajustando también durante el entrenamiento:
\begin{equation}
    \mathbf{Y} = \text{Concatenar} [\mathbf{A}_1,\mathbf{A}_2, \dots, \mathbf{A}_H]\mathbf{\Omega_o}
    \label{multiples_capas}
\end{equation}
donde la matriz resultante de concatenar las matrices $\mathbf{A}_h$ 
tiene dimensión \mbox{$N \times H \cdot D_i$} y, en consecuencia, $\mathbf{\Omega_o}$ tiene dimensión $H \cdot D_i \times D_i$ conservando así la dimensión $N \times D_i$ de X.

En la práctica,  cada una de estas capas de autoatención $\mathbf{A}_h$, 
tiene una dimensión inicial $N \times D_a$ donde $D_a = \frac{D_i}{H}$ con 
el objetivo de que la matriz concatenada tenga dimensión $N \times D_i$ 
(ver figura \ref{fig:concatenar_matrices}). Además, de esta manera el coste 
computacional total de evaluar las $H$ capas en paralelo se mantiene 
aproximadamente igual al coste que tendría evaluar una única capa de atención.
Esto evita que el coste de las operaciones matriciales se dispare al utilizar 
múltiples capas.

\begin{figure}[t]
    \vspace{0.8cm}
    \centering
    \FlujoAtencionMultiple
    \caption{Flujo de datos en la atención múltiple. }
    \label{fig:Flujo_atencion_multiple}
    \vspace{0.8cm}
\end{figure}




\section{Transformador: Arquitectura}

La atención múltiple definida en este último apartado es el eje central de la arquitectura en un transformador. Sin embargo, es bien sabido que las redes neuronales suelen mejorar su desempeño a medida que aumentan en profundidad, es decir, en número de capas, y los transformadores no son una excepción.

Para mejorar el entrenamiento y evitar que el gradiente se desvanezca en la profundidad de la red, implementamos una \textit{conexión residual}: una conexión que permite que la información de una capa anterior se salte las operaciones intermedias para sumarse directamente a la salida de estas. Esta suma se puede hacer gracias a que las \textit{capas de autoatención} no modifican la dimensión de la entrada \textbf{X}. Después se aplica una normalización de capas que ayuda a estabilizar y acelerar el aprendizaje:
\begin{equation}
    \mathbf{Z} = \text{LayerNorm}[\mathbf{Y}[\mathbf{X}]+\mathbf{X}]
\end{equation}

Posteriormente, y con el objetivo de aumentar la no-linealidad del modelo, pasamos \textbf{Z} por una red neuronal $ \y = f[\x,\phi]$ donde $\x,\y \in \mathbb{R}^{D_i}$:
\begin{equation}
    \mathbf{\Tilde{X}} = \text{LayerNorm}[f[\mathbf{Z},\phi] + \mathbf{Z}]
    \label{transformer_ecuacion}
\end{equation}
Y esta sería la forma final de un transformador. 

\begin{figure}[t]
    \centering
    \DiagramaTransformer
    \caption{Arquitectura del transformador.}
    \label{fig:diagrama_transformer}
\end{figure}

Habitualmente, la arquitectura de estos modelos no se compone de un 
único bloque, sino que se construye concatenando múltiples transformadores 
apilados, donde la salida de uno constituye la entrada del siguiente. 
Desde su introducción original \cite{vaswani2017attention_arxiv}, el modelo 
fue diseñado expresamente apilando $6$ transformadores idénticoss, pero 
aunque estas capas comparten exactamente la misma estructura, cada capa utiliza parámetros, pesos y sesgos completamente independientes que el modelo optimiza de manera individual durante el entrenamiento.

El hecho de que cada capa aprenda parámetros distintos permite que cada nivel se ``especialice'' en un grado de abstracción diferente conforme la información avanza por el modelo. Las primeras capas tienden a enfocarse en relaciones locales y patrones sintácticos básicos, mientras que las capas más profundas utilizan esa base para construir representaciones semánticas globales y complejas. Si todas las capas compartieran los mismos parámetros el modelo se limitaría a aplicar siempre la misma transformación, lo que limitaría  su capacidad de aprendizaje.

El éxito de esta  arquitectura se determinó y verificó de manera empírica. 
Ya en el artículo original los autores demostraron mediante diversas 
pruebas que los modelos más grandes y profundos ofrecían una mejor calidad 
en las tareas de traducción. En la actualidad, esta observación ha 
evolucionado hasta consolidarse en la denominada hipótesis del escalado 
(\textit{scaling hypothesis}) \cite{kaplan2020scaling}. Esta hipótesis 
postula que, al aumentar a gran escala el número de parámetros y capas 
del modelo, junto a bases de datos enormes, se logran incrementos  en el 
rendimiento sin necesidad de alterar la arquitectura básica del transformador.
Como evidencia empírica contemporánea de este fenómeno, modelos de lenguaje de gran tamaño como GPT-3 han llevado este concepto al extremo, utilizando una pila de 96 transformadores.





